\section{La retraite elle-m\^eme}\label{sec:retraite}
% Chapitre 8 (tache 10, passage chapitre 08). Figures: fig_rente_alter, fig_bruecke,
% fig_bruecke_anticipee (nouvelle: depart anticipe au taux d'epargne BRUECKE_QUOTE_NOM),
% fig_puffer (MEME calcul Monte Carlo que \ErfolgsquoteBasis/\ErfolgsquoteMitPuffer du
% chapitre 7), fig_regle_quatre_histoire (nouvelle: test historique du retrait de
% TAUX_RETRAIT). Tableau tab:retrait-contre-dividendes ecrit ici. Aucun chiffre tape:
% macros de data/kennzahlen.tex (kapitel_8 et _macros_redaction_kapitel_8 dans
% scripts/rechnung_retraite.py). La comparaison du tableau porte sur les MEMES annees de
% depart en retraite que le rejeu du chapitre 7 (cache _RUECKSPIEL, aucun second rejeu).

Ce chapitre r\'epond \`a la question~: une fois l'objectif atteint, de quoi vis-tu jusqu'\`a
l'\^age l\'egal, que t'apporte ensuite la rente l\'egale, quelle r\'eserve faut-il garder et selon
quelle r\`egle puiser dans ton capital~? Le chapitre~\ref{sec:risques} s'est conclu sur une r\'eserve
de liquidit\'es qui prot\`ege ton revenu dans la plupart des baisses du dividende, et sur un mod\`ele
qui mesure la solidit\'e de ce dividende, pas la date de ton d\'epart. Ce chapitre revient \`a ton
calendrier. La section~\ref{sec:retraite-rente} estime ta rente l\'egale selon l'\^age auquel tu
cesses de travailler, et la section~\ref{sec:retraite-pont} en d\'eduit la phase pont, les ann\'ees qui
s\'eparent ton d\'epart du versement de cette rente. La section~\ref{sec:retraite-reserve} chiffre la
taille de la r\'eserve, la section~\ref{sec:retraite-quatre} soumet la r\`egle des \TauxRetrait~\% \`a
l'histoire comme le chapitre~\ref{sec:risques} l'a fait pour les dividendes, et la
section~\ref{sec:retraite-regles} en tire tes r\`egles de retrait.

\subsection{La rente l\'egale selon l'\^age o\`u tu arr\^etes}\label{sec:retraite-rente}

Tout ce plan vise un revenu tir\'e de ton capital. Tu cotises pourtant \`a l'assurance pension
l\'egale d\`es ton premier emploi, et cette rente te sera vers\'ee \`a partir de l'\^age l\'egal, que
tu travailles encore ou non \`a ce moment-l\`a.

\annahme{Rente l\'egale estim\'ee par points~: chaque ann\'ee de travail rapporte ton salaire brut
divis\'e par le salaire moyen de 2026, \Durchschnittsentgelt~€\QH{durchschnittsentgelt_2026}, et
chaque point vaut \Rentenwert~€ de rente mensuelle\QH{rentenwert_2026}. Ton salaire ne compte que
jusqu'au plafond de cotisation de l'assurance pension, \BbgRvJahr~€ par an\QH{sv_arbeitnehmer}, que
ton salaire atteint \`a \AgePlafondRv~ans. Salaire du chapitre~\ref{sec:accumulation}, carri\`ere
continue depuis \AgeDebutCotisation~ans, sans ch\^omage ni interruption. Ces valeurs de 2026 sont
tenues constantes en euros de 2026, comme les seuils de l'assurance maladie~: le plan suppose que la
valeur du point garde son pouvoir d'achat, ce que rien ne garantit dans une population qui
vieillit. Rente vers\'ee \`a partir de l'\^age l\'egal de \AgeLegal~ans\QH{regelaltersgrenze}, sans
abattement~; le plan ne mod\'elise aucun d\'epart anticip\'e de la rente elle-m\^eme.}

L'\'equation~\eqref{eq:rente-points-legale} r\'esume ce calcul pour un arr\^et du travail \`a
l'\^age~$a$~; elle additionne les points de chaque ann\'ee travaill\'ee.
\begin{equation}
  R(a) = V \sum_{t=a_0}^{a-1} \frac{\min(S_t,\,P)}{D} \text{,} \label{eq:rente-points-legale}
\end{equation}
o\`u $R(a)$ est la rente brute mensuelle vers\'ee \`a partir de \AgeLegal~ans, $a_0$ l'\^age de ton
premier emploi, \AgeDebutCotisation~ans, $S_t$ ton salaire brut annuel \`a l'\^age~$t$, $P$ le
plafond de cotisation, $D$ le salaire moyen et $V$ la valeur du point. Ta premi\`ere ann\'ee de
travail rapporte \PointsEntree~point, une ann\'ee au plafond \PointsPlafond~point~; aucune ann\'ee ne
peut rapporter davantage, quel que soit ton salaire. La figure~\ref{fig:rente-selon-arret} applique
l'\'equation \`a chaque \^age d'arr\^et de \RenteAgeMinGraphique{} \`a \AgeLegal~ans.

\linienfigur{fig_rente_alter}[width=0.95\linewidth, xlabel={\^Age auquel tu arr\^etes de travailler (ans)},
  ylabel={Rente l\'egale brute (euros de 2026 par mois)}, scaled y ticks=false, ymin=0,
  legend pos=north west, legend cell align=left, legend style={fill=white}]%
  {ausstiegsalter}{rente_monat}%
  {{Rente brute vers\'ee d\`es \AgeLegal~ans}}%
  {Rente l\'egale mensuelle brute vers\'ee \`a partir de \AgeLegal~ans selon l'\^age auquel tu arr\^etes
  de travailler (carri\`ere continue depuis \AgeDebutCotisation~ans, valeurs de 2026, euros de
  2026).\label{fig:rente-selon-arret}}

Ce graphique montre qu'en travaillant jusqu'\`a \AgeLegal~ans, tu toucherais \RenteBeiZielalter~€
bruts par mois, soit \RenteLegalePartCible~\% de ta cible de \ZielMonat~€. Un arr\^et \`a 60~ans la
ram\`ene \`a \RenteArretSechzig~€, \`a 50~ans \`a \RenteArretFuenfzig~€ et \`a 40~ans \`a
\RenteArretVierzig~€. La pente de la courbe augmente avec ton salaire jusqu'\`a \AgePlafondRv~ans,
puis reste constante~: au-del\`a du plafond, chaque ann\'ee de travail ajoute la m\^eme rente,
\RenteDerniereAnnee~€ par mois pour la derni\`ere. \bk{Chaque ann\'ee de retraite anticip\'ee te
co\^ute donc de la rente \`a vie, en plus des ann\'ees de pont \`a financer.}

Ces montants sont bruts. La rente est soumise \`a l'imp\^ot sur le revenu et aux cotisations
d'assurance maladie des retrait\'es, que ce plan ne calcule pas~; elle n'est donc pas directement
comparable aux \ZielMonat~€ nets de ta cible, mais elle en couvre une grande partie d\`es que ta
carri\`ere est compl\`ete.

\subsection{La phase pont jusqu'\`a l'\^age l\'egal}\label{sec:retraite-pont}

La phase pont est la p\'eriode entre le jour o\`u tu cesses de travailler et celui o\`u ta rente
l\'egale commence~: pendant ces ann\'ees, ton capital doit te faire vivre seul.

\annahme{Pendant le pont, ton revenu est le dividende net de l'ann\'ee o\`u tu atteins l'objectif,
qui garde ensuite sa valeur r\'eelle, comme au chapitre~\ref{sec:capital}. Tu es alors assur\'e
volontaire, et ce revenu est d\'ej\`a net de la cotisation maladie du chapitre~\ref{sec:fiscalite}. Le
plan garde les m\^emes r\`egles apr\`es \AgeLegal~ans~; l'assurance maladie des retrait\'es ob\'eit \`a
d'autres r\`egles, qu'il n'int\`egre pas et que tu devras faire v\'erifier.}

Dans le sc\'enario de base, le pont dure \BrueckeJahre~ann\'ee~: tu atteins l'objectif \`a
\AlterYannBasis~ans, qui est l'\^age l\'egal. La figure~\ref{fig:pont-scenario-base} montre les deux
revenus qui en r\'esultent.

\linienfigur{fig_bruecke}[width=0.95\linewidth, xlabel={Ton \^age (ans)}, ylabel={Revenu mensuel (euros de 2026)},
  scaled y ticks=false, ymin=0, enlarge y limits={upper, value=0.2},
  cycle list={{blue, thick},{red, thick, dashed}},
  legend pos=south east, legend cell align=left, legend style={fill=white}]%
  {alter}{dividende,rente}%
  {{Dividendes nets},{Rente l\'egale brute}}%
  {Revenu mensuel du sc\'enario de base \`a partir de l'objectif, atteint \`a \AlterYannBasis~ans~:
  dividendes nets et rente l\'egale brute, sans phase pont (r\'epartition \ReferenzStrategie,
  \ReferenzSparquote~\% d'\'epargne, euros de 2026).\label{fig:pont-scenario-base}}

Ce graphique montre deux revenus qui commencent la m\^eme ann\'ee, \`a \AlterYannBasis~ans~:
\RevenuNetYannBasis~€ nets de dividendes et \RenteBeiZielalter~€ bruts de rente l\'egale. \bk{Le
sc\'enario de base fait donc travailler ton capital pour un revenu que ta rente couvre d\'ej\`a en
grande partie.} La cible de ce plan ne compte pas la rente~: le capital du chapitre~\ref{sec:capital}
est dimensionn\'e comme si elle n'existait pas. Ce choix est prudent, parce qu'il te rend
ind\'ependant du niveau futur des rentes, mais il a un prix~: pass\'e \AgeLegal~ans, tu disposerais de
bien plus que ta cible, et une partie de l'\'epargne de tes ann\'ees de travail n'aurait servi qu'\`a
ce surplus. Un plan qui compterait la rente demanderait \`a ton capital de couvrir toute la cible
jusqu'\`a l'\^age l\'egal, puis seulement l'\'ecart avec la rente nette~; ce document ne le calcule
pas, et c'est une piste \`a chiffrer avec un conseiller.

Le pont n'existe que si tu pars plus t\^ot. Au chapitre~\ref{sec:accumulation}, un taux d'\'epargne de
\SparquoteFuenfzig~\% avec la r\'epartition \ReferenzStrategie{} atteint l'objectif \`a
\AlterMaisonSiebzigFuenfzig~ans, en \BrueckeAnneeDebutFuenfzig, avec un capital de
\KapitalBrueckeFuenfzig~€. La figure~\ref{fig:pont-depart-anticipe} montre le revenu qui en r\'esulte.

\linienfigur{fig_bruecke_anticipee}[xlabel={Ton \^age (ans)}, ylabel={Revenu mensuel (euros de 2026)},
  scaled y ticks=false, ymin=0, enlarge y limits={upper, value=0.25},
  cycle list={{blue, thick},{red, thick, dashed},{black, dotted}},
  every axis plot post/.append style={const plot},
  legend pos=south east, legend cell align=left, legend style={fill=white}]%
  {alter}{dividende,rente,cible}%
  {{Dividendes nets},{Rente l\'egale brute},{Cible de \ZielMonat~€}}%
  {Revenu mensuel apr\`es un d\'epart anticip\'e \`a \AlterMaisonSiebzigFuenfzig~ans~: dividendes nets
  seuls pendant le pont, puis rente l\'egale brute \`a partir de \AgeLegal~ans (r\'epartition
  \ReferenzStrategie, \SparquoteFuenfzig~\% d'\'epargne, euros de 2026).\label{fig:pont-depart-anticipe}}

Ce graphique montre un pont de \BrueckeJahreFuenfzig~ans, de \AlterMaisonSiebzigFuenfzig{} \`a
\AgeLegal~ans, pendant lequel tes dividendes, \RevenuNetBrueckeFuenfzig~€ nets par mois, sont ton
seul revenu et d\'epassent \`a peine la cible. La rente arrive en \BrueckeAnneeRente, mais elle ne
vaut que \RenteBrueckeFuenfzig~€ bruts, \RenteBrueckePerteMois~€ de moins que si tu avais travaill\'e
jusqu'\`a \AgeLegal~ans, soit \RenteBrueckePerteProzent~\% de moins, et cela pour le reste de ta vie.
Ici, le pont n'exige aucun capital suppl\'ementaire, puisque les dividendes couvrent seuls la cible
d\`es le d\'epart~; son co\^ut est cette rente r\'eduite. Son risque est ailleurs~: une baisse du
dividende dans les premi\`eres ann\'ees de retraite, la plus dangereuse selon la
section~\ref{sec:risques-sequence}, tomberait pendant le pont, quand tu n'as aucun autre revenu. Au
taux d'\'epargne minimal qui permet de partir avant l'\^age l\'egal, \QuoteAnticipeeMaisonSiebzig~\%,
tu partirais \`a \AlterAnticipeeMaisonSiebzig~ans~: le pont ne durerait que
\BrueckeJahreAnticipee~ans, et ta rente vaudrait \RenteBrueckeAnticipee~€ bruts.

\subsection{Quelle r\'eserve de liquidit\'es~?}\label{sec:retraite-reserve}

La section~\ref{sec:risques-reserve} a montr\'e qu'une r\'eserve de \PufferJahreMit~ans de d\'epenses
fait passer la r\'eussite du mod\`ele sans imp\^ot de \ErfolgsquoteBasis~\% \`a
\ErfolgsquoteMitPuffer~\%. Reste \`a savoir quelle taille lui donner. La
figure~\ref{fig:reussite-selon-reserve} la fait varier de z\'ero \`a \PufferJahreMax~ans, sur les m\^emes
trajectoires et avec les m\^emes r\`egles que le Monte Carlo du chapitre~\ref{sec:risques}
(tableau~\ref{tab:projection-contre-simulation}).

\annahme{R\'eserve mesur\'ee en ann\'ees de cible, \ZielJahr~€ chacune, constitu\'ee au d\'epart en
retraite et jamais reconstitu\'ee. Elle ne comble que ce qui manque pour remonter au seuil de
\SeuilReussite~\% de la cible. Le mod\`ele lui garde sa valeur r\'eelle, ce qui suppose un placement
dont l'int\'er\^et compense au moins l'inflation~: un compte non r\'emun\'er\'e perdrait chaque ann\'ee
du pouvoir d'achat. Mod\`ele sans imp\^ot, march\'e am\'ericain seul, retraite de \McRenteAns~ans.}

\linienfigur{fig_puffer}[xlabel={R\'eserve de liquidit\'es (ann\'ees de d\'epenses)},
  ylabel={Probabilit\'e de r\'eussite (\%)}, y filter/.expression={y*1e2}, xtick=data, ymin=0,
  enlarge y limits={upper, value=0.15}, enlarge x limits=0.1,
  every axis plot post/.append style={mark=*, mark size=1.8pt},
  nodes near coords, every node near coord/.append style={font=\scriptsize, anchor=south east,
    /pgf/number format/fixed, /pgf/number format/fixed zerofill, /pgf/number format/precision=1},
  legend pos=south east, legend cell align=left, legend style={fill=white}]%
  {puffer_jahre}{erfolgsquote}%
  {{R\'eussite sur \McRenteAns~ans}}%
  {Probabilit\'e que le revenu reste au-dessus de \SeuilReussite~\% de la cible pendant \McRenteAns~ans
  de retraite, selon la taille de la r\'eserve de liquidit\'es (Monte Carlo sans imp\^ot du
  chapitre~\ref{sec:risques}, \McTrajectoires{} trajectoires).\label{fig:reussite-selon-reserve}}

Ce graphique montre des gains d\'ecroissants. La premi\`ere ann\'ee de r\'eserve ajoute
\PufferGainEins~points de r\'eussite, de \ErfolgsquoteBasis{} \`a \ErfolgsquotePufferEins~\%, la
deuxi\`eme \PufferGainMit~points, jusqu'\`a \ErfolgsquoteMitPuffer~\%, et la troisi\`eme
\PufferGainMax~points, jusqu'\`a \ErfolgsquotePufferMax~\%. Chaque ann\'ee co\^ute pourtant autant \`a
\'epargner, \ZielJahr~€, soit \PufferMontantMax~€ pour \PufferJahreMax~ans, en plus du capital du
chapitre~\ref{sec:capital}. \bk{Une ann\'ee de d\'epenses en r\'eserve apporte l'essentiel de la
protection~; au-del\`a, chaque euro mis de c\^ot\'e prot\`ege de moins en moins.} Ces taux sont ceux du
mod\`ele sans imp\^ot, qui ne dit rien de ta date de d\'epart~; c'est l'\'ecart entre eux qui compte.
La r\'eserve sert surtout pendant le pont~: \`a partir de \AgeLegal~ans, ta rente, qui ne d\'epend
d'aucun march\'e, joue en partie le m\^eme r\^ole.

\subsection{La r\`egle des \TauxRetrait~\% \`a l'\'epreuve de l'histoire}\label{sec:retraite-quatre}

Le chapitre~\ref{sec:capital} a oppos\'e deux fa\c{c}ons de vivre de ton capital, et le
chapitre~\ref{sec:risques} n'a rejou\'e que la premi\`ere. Cette section rejoue la seconde~: un ETF
actions dont tu retires chaque ann\'ee \TauxRetrait~\% du capital de d\'epart, un montant ensuite
constant en pouvoir d'achat, quels que soient les march\'es.

\annahme{Donn\'ees de Robert Shiller\QH{shiller_url}~: rendement r\'eel total de
l'\IndiceActionsUs{}, dividendes r\'einvestis, jusqu'en \ShillerFin. Pour chaque ann\'ee de d\'epart,
tu retires en d\'ebut d'ann\'ee \TauxRetrait~\% du capital initial, ajust\'es de l'inflation, puis le
capital restant re\c{c}oit le rendement r\'eel de l'ann\'ee. Le capital tient s'il couvre tous les
retraits pendant \RetraitHistAnsDreissig{} ou \RetraitHistAnsVierzig~ans~; d\'eparts de
\RetraitHistPremierDreissig{} \`a \RetraitHistDernierDreissig{} pour \RetraitHistAnsDreissig~ans, et
jusqu'en \RetraitHistDernierVierzig{} pour \RetraitHistAnsVierzig~ans. L'imp\^ot sur les parts vendues
est pay\'e sur le retrait lui-m\^eme, comme au chapitre~\ref{sec:capital}. \textbf{Limites
d\'eclar\'ees~: march\'e am\'ericain seul, portefeuille enti\`erement en actions, sans frais ni imp\^ot
pr\'elev\'e sur le capital.} Les frais d'un ETF et l'imp\^ot annuel de la Vorabpauschale retirent
chaque ann\'ee une petite part du capital, ce qui ne peut qu'abaisser les taux qui suivent.}

\linienfigur{fig_regle_quatre_histoire}[xlabel={Ann\'ee de d\'epart en retraite},
  ylabel={Capital r\'eel final (multiple du capital initial)},
  xticklabel style={/pgf/number format/set thousands separator={}},
  table/empty cells with={nan}, unbounded coords=discard, ymin=0,
  cycle list={{blue},{red}}, every axis plot post/.append style={only marks, mark=*,
    mark size=1.1pt},
  legend cell align=left, legend style={at={(0.5,-0.22)}, anchor=north, legend columns=2,
    /tikz/every even column/.append style={column sep=0.5cm}}]%
  {start}{fin_dreissig,fin_vierzig}%
  {{Apr\`es \RetraitHistAnsDreissig~ans de retraits},{Apr\`es \RetraitHistAnsVierzig~ans de retraits}}%
  {Capital r\'eel restant apr\`es \RetraitHistAnsDreissig{} et \RetraitHistAnsVierzig~ans de retraits de
  \TauxRetrait~\% du capital initial, selon l'ann\'ee de d\'epart~; z\'ero signifie un capital
  \'epuis\'e avant la fin (indice \IndiceActionsUs, sans frais, march\'e am\'ericain
  seul).\label{fig:regle-retrait-histoire}}

% Phrases liees aux donnees actuelles (a relire si la serie change): "autour de 1929",
% "fin des annees 1960 et debut des annees 1970", "aucun depart n'echoue" au taux prudent.
Ce graphique montre que, sur \RetraitHistAnsDreissig~ans, le capital tient pour
\RetraitHistTenusDreissig{} des \RetraitHistDebutsDreissig{} ann\'ees de d\'epart, soit
\RetraitHistReussiteDreissig~\%, et s'\'epuise pour les d\'eparts de \RetraitHistEchecsListeDreissig.
Sur \RetraitHistAnsVierzig~ans, il tient pour \RetraitHistTenusVierzig{} d\'eparts sur
\RetraitHistDebutsVierzig, soit \RetraitHistReussiteVierzig~\%, et s'\'epuise pour ceux de
\RetraitHistEchecsListeVierzig. Le pire d\'epart est \RetraitHistPireDebutVierzig~: le capital n'a
couvert que \RetraitHistPireDureeVierzig~ann\'ees de retraits. Les \'echecs se concentrent sur des
retraites commenc\'ees juste avant une longue p\'eriode de rendements r\'eels faibles, autour de 1929
et surtout \`a la fin des ann\'ees 1960 et au d\'ebut des ann\'ees 1970, au seuil de la d\'ecennie
d'inflation (section~\ref{sec:risques-baisses}).

Dans le cas m\'edian, au contraire, le capital r\'eel final vaut \RetraitHistFinMedianeFoisDreissig{}
fois le capital de d\'epart apr\`es \RetraitHistAnsDreissig~ans, et \RetraitHistFinMedianeFoisVierzig{}
fois apr\`es \RetraitHistAnsVierzig~ans~: la r\`egle est pens\'ee pour les pires s\'equences, et dans
la plupart des histoires elle laisse un gros capital inutilis\'e. Au taux prudent de
\TauxRetraitPrudent~\% du chapitre~\ref{sec:capital}, aucun d\'epart n'\'echoue~: la r\'eussite vaut
\RetraitHistReussitePrudentVierzig~\% sur \RetraitHistAnsVierzig~ans, et dans le pire cas il reste
encore \RetraitHistPrudentFinMinVierzig~\% du capital initial. Ce taux demande cependant
\KapitalEtfPrudent~€ de capital au lieu de \KapitalEtfReferenz~€.

Ce test ne r\'epond pas \`a la m\^eme question que le rejeu des dividendes de la
section~\ref{sec:risques-rejeu}. La r\`egle des \TauxRetrait~\% ne baisse jamais ton revenu~; elle
\'echoue quand le capital est vide, et cet \'echec est d\'efinitif. La strat\'egie des dividendes ne
vend jamais rien~; elle \'echoue d\`es que son revenu passe une seule ann\'ee sous \SeuilReussite~\% de
l'objectif, mais le capital reste, et le revenu peut remonter. Le
tableau~\ref{tab:retrait-contre-dividendes} compare pourtant les deux issues sur les m\^emes
\RetraitHistMemesDepartsJuges{} d\'eparts jug\'es que le rejeu, avec une retraite de
\RetraitHistAnsVierzig~ans qui commence l'ann\'ee suivant l'objectif.

\begin{table}[!tb]
  \centering\small
  \caption{Issue des \RetraitHistMemesDepartsJuges{} d\'eparts jug\'es du rejeu du
  chapitre~\ref{sec:risques}~: revenu de dividendes contre r\`egle des \TauxRetrait~\% sur les m\^emes
  retraites de \RetraitHistAnsVierzig~ans (indice \IndiceActionsUs, sans imp\^ot, nombre de
  d\'eparts).}\label{tab:retrait-contre-dividendes}
  \begin{tabularx}{\linewidth}{>{\bfseries}l *{3}{>{\centering\arraybackslash}X}}
    \toprule
    \rowcolor[gray]{0.9}
    \textbf{Revenu de dividendes} & \textbf{R\`egle des \TauxRetrait~\%~: le capital tient}
      & \textbf{R\`egle des \TauxRetrait~\%~: capital \'epuis\'e} & \textbf{Total} \\
    \midrule
    Tenu toute la retraite & \RetraitHistLesDeuxTiennent & \RetraitHistSeulDividendeTient
      & \RueckspielTenus \\
    Pass\'e sous le seuil & \RetraitHistSeulRetraitTient & \RetraitHistAucunNeTient & \RueckspielEchecs \\
    Total & \RetraitHistMemesDepartsTenus & \RetraitHistMemesDepartsEchecs & \RetraitHistMemesDepartsJuges \\
    \bottomrule
  \end{tabularx}
\end{table}

Ce tableau montre que, sur ces d\'eparts, la r\`egle des \TauxRetrait~\% tient dans
\RetraitHistMemesDepartsTenus{} cas, soit \RetraitHistMemesDepartsProzent~\%, contre
\RueckspielTenus{} cas, \RueckspielTenusProzent~\%, pour le revenu de dividendes. Les deux
\'echouent ensemble pour les retraites commenc\'ees en \RetraitHistAnneesAucun. Le revenu de
dividendes tient seul pour celles de \RetraitHistAnneesSeulDividende, quand les rendements r\'eels du
march\'e ont \'et\'e trop faibles pour le capital sans que le dividende passe sous le seuil. La r\`egle
tient seule pour celles de \RetraitHistAnneesSeulRetrait, o\`u le dividende est pass\'e sous le seuil
sans que le capital s'\'epuise. \bk{Sur ce test historique, la r\`egle des \TauxRetrait~\% ne fait pas
moins bien que les dividendes seuls~: elle tient plus souvent, avec moins de capital
(chapitre~\ref{sec:capital}).} La presse sp\'ecialis\'ee pr\'esente parfois la strat\'egie des
dividendes comme celle qui r\'eduit au minimum le risque de s\'equence et de
faillite\QP{justetf_entnahmestrategien}~; c'est vrai pour la faillite, puisque le capital n'est jamais
vendu, mais pas pour le revenu, que la s\'equence des rendements frappe \`a travers le dividende
lui-m\^eme.

Deux prudences s'imposent. Les \RetraitHistMemesDepartsJuges{} d\'eparts ne correspondent qu'\`a
\RetraitHistMemesDepartsAnnees{} ann\'ees de retraite distinctes, qui se chevauchent largement~: ces
pourcentages d\'ecrivent une poign\'ee d'histoires, pas des dizaines d'exp\'eriences ind\'ependantes.
Et l'\'echec n'a pas la m\^eme gravit\'e des deux c\^ot\'es~: un revenu de dividendes pass\'e sous le
seuil reste un revenu, avec tout le capital derri\`ere lui, alors qu'un capital \'epuis\'e ne verse plus
rien. Les deux tests partagent enfin les m\^emes limites~: march\'e am\'ericain seul, sans frais, et
sans la fiscalit\'e allemande, qui abaisse le taux de retrait effectif\QP{getmad_fire}.

\subsection{Tes r\`egles de retrait}\label{sec:retraite-regles}

Ces r\'esultats se traduisent en quelques r\`egles, \`a fixer avant le d\'epart pour ne pas avoir \`a
les inventer pendant une crise.

Si tu suis la strat\'egie des dividendes, tu vis des dividendes nets de l'ann\'ee et tu ne vends
aucune part. Quand le dividende baisse, tu acceptes de r\'eduire tes d\'epenses jusqu'\`a
\SeuilReussite~\% de ta cible, et tu ne puises dans la r\'eserve que pour ne pas descendre plus
bas~: c'est la r\`egle que simule le chapitre~\ref{sec:risques}, et la seule dont ce document mesure
l'effet. Le mod\`ele ne reconstitue jamais la r\'eserve~; la remplir de nouveau les bonnes ann\'ees
est une pr\'ecaution de bon sens, que le calcul ne chiffre pas.

Si tu suis la r\'ef\'erence en ETF, tu retires chaque ann\'ee le m\^eme montant r\'eel,
\TauxRetrait~\% du capital de d\'epart ajust\'es de l'inflation, sans le relever apr\`es une bonne
ann\'ee ni le couper apr\`es une mauvaise~: c'est la r\`egle test\'ee dans la
section~\ref{sec:retraite-quatre}. Pour \ZielMonat~€ nets par mois, ce retrait vaut
\RetraitBrutEtf~€ bruts par an, imp\^ot et cotisation maladie compris (chapitre~\ref{sec:capital}).
Si une crise survient dans les premi\`eres ann\'ees, la tentation de tout vendre ou d'abandonner la
r\`egle sera forte~; sur l'histoire am\'ericaine, la tenir a suffi dans \RetraitHistReussiteVierzig~\%
des d\'eparts sur \RetraitHistAnsVierzig~ans.

Dans les deux cas, ta rente l\'egale change la donne \`a \AgeLegal~ans~: elle couvre alors une grande
partie de ta cible et ne d\'epend d'aucun march\'e. Tu pourras r\'eduire ce que tu tires de ton
capital, ou garder ce revenu comme marge contre une baisse future du niveau des rentes. Le plan ne
chiffre pas cette seconde phase~; il faudra la recalculer quelques ann\'ees avant l'\^age l\'egal,
avec le relev\'e de carri\`ere que t'enverra l'assurance pension.

\bk{Retiens de ce chapitre qu'au sc\'enario de base il n'y a pas de pont, puisque l'objectif tombe \`a
l'\^age l\'egal, et que ta rente l\'egale, \RenteBeiZielalter~€ bruts par mois, s'ajoute alors \`a des
dividendes d\'ej\`a suffisants. Un d\'epart \`a \AlterMaisonSiebzigFuenfzig~ans demande
\SparquoteFuenfzig~\% d'\'epargne et r\'eduit cette rente de \RenteBrueckePerteProzent~\% \`a vie.}
Une ann\'ee de d\'epenses en r\'eserve apporte l'essentiel de la protection mesurable. Enfin, rejou\'ee
sur l'histoire am\'ericaine, la r\`egle des \TauxRetrait~\% tient \RetraitHistAnsVierzig~ans dans
\RetraitHistReussiteVierzig~\% des d\'eparts~; sur les \RetraitHistMemesDepartsJuges{} d\'eparts du
rejeu des dividendes, elle tient dans \RetraitHistMemesDepartsProzent~\% des cas, contre
\RueckspielTenusProzent~\% pour le revenu de dividendes, et avec moins de capital. Le chapitre~\ref{sec:presse} confronte ces r\'esultats \`a ce qu'\'ecrivent
les magazines et la recherche.

\Yannbox{\AnneeYannBasis}{En \AnneeYannBasis, \`a \AlterYannBasis~ans, Yann arr\^ete de travailler
l'ann\'ee m\^eme o\`u sa rente l\'egale commence~: pas de pont. Ses dividendes lui versent
\RevenuNetYannBasis~€ nets par mois et sa rente \RenteBeiZielalter~€ bruts. S'il avait \'epargn\'e
\SparquoteFuenfzig~\% de son salaire, il serait parti en \BrueckeAnneeDebutFuenfzig, \`a
\AlterMaisonSiebzigFuenfzig~ans, aurait v\'ecu \BrueckeJahreFuenfzig~ans de ses seuls dividendes,
\RevenuNetBrueckeFuenfzig~€ nets par mois, puis touch\'e \RenteBrueckeFuenfzig~€ bruts de rente, soit
\RenteBrueckePerteMois~€ de moins. Une r\'eserve d'un an, \ZielJahr~€, porterait la r\'eussite du
mod\`ele sans imp\^ot de \ErfolgsquoteBasis~\% \`a \ErfolgsquotePufferEins~\%. Avec un ETF monde et un
retrait de \TauxRetrait~\%, son capital aurait tenu \RetraitHistAnsVierzig~ans dans
\RetraitHistReussiteVierzig~\% des d\'eparts historiques de l'\IndiceActionsUs.}
